\section{Conclusion}
\label{sec:conclusion}

We studied CAP's contrastive pruning objective \citep{xu2022cap} --- the PrC / FiC
/ SnC decomposition, which is CAP's contribution and not ours --- outside the
setting it was developed in: on convolutional vision backbones rather than
language models, under three different pruning criteria rather than one, and at
sparsities up to 99.9\% rather than 97\%. The comparison was made against an
iterative pruning baseline sharing the same initialisation, schedule, criterion,
epoch budget and fine-tune, differing only in the training objective, so that the
measured difference is attributable to the objective alone.

The evidence is bounded in ways we have tried to state precisely rather than
minimise: one dataset, a compute budget well below the published comparators',
seed replication on the headline cells only, and a trainable projection head that
makes the accuracy result defensible but the representation-preservation reading
of the loss not. Within those bounds, the useful output of this work is a
controlled measurement of what changing the retraining objective is worth as a
network is pruned past the point where cross-entropy alone can rebuild it.
